
我们指出，式~\eqref{eq:7-4-35}
与~\eqref{eq:3-8-33}
几乎相同。
唯一的差别在于，
第一项多出符号
$(-1)^{\deg\mathbf y}$，
而第二项多出符号
$(-1)^{
\deg\mathbf y_{c_1}^{(2;1)}
\deg\mathbf y_{c_1}^{(2;2)}}$。
为消去它们，定义
\begin{equation}
\mathfrak q_{\ell,k}
(y_1\otimes\cdots\otimes y_\ell)
(\mathbf x)
=\mathfrak q^{\circ}_{\ell,k}
(y_\ell\otimes\cdots\otimes y_1)
(\mathbf x),
\tag{7.4.39.1}\label{eq:7-4-39-1}
\end{equation}
以此代替 $\mathfrak q^{\circ}_{\ell,k}$\liftIndex{I289}；
并以
\begin{equation}
\widehat\delta
(y_1\otimes\cdots\otimes y_k)
=\sum_{i=1}^{k}
(-1)^{\deg y_1+\cdots+\deg y_{i-1}}
y_1\otimes\cdots\otimes
\delta y_i\otimes\cdots\otimes y_k
\tag{7.4.39.2}\label{eq:7-4-39-2}
\end{equation}
代替 $\widehat\delta^{\circ}$。
事实上，若记
\[
\operatorname{Op}
(y_1\otimes\cdots\otimes y_\ell)
=y_\ell\otimes\cdots\otimes y_1,
\]
则
\begin{equation}
\operatorname{Op}
\bigl(\widehat\delta^{\circ}(\mathbf y)\bigr)
=(-1)^{\deg\mathbf y}
\widehat\delta
\bigl(\operatorname{Op}(\mathbf y)\bigr).
\tag{7.4.39.3}\label{eq:7-4-39-3}
\end{equation}
因此，在式~\eqref{eq:7-4-35}
的第一项中令
$\mathbf y=\operatorname{Op}(\mathbf y')$，
它便成为
\[
\mathfrak q_{\ell,k}
\bigl(\widehat\delta(\mathbf y')
\otimes\mathbf x\bigr).
\]

另一方面，有
\[
\begin{aligned}
\Delta\operatorname{Op}(\mathbf y)
&=\sum_{c_1}
\operatorname{Op}
(\mathbf y_{c_1}^{(2;2)})
\otimes
\operatorname{Op}
(\mathbf y_{c_1}^{(2;1)})\\
&=\sum_{c_1}
(-1)^{
\deg\mathbf y_{c_1}^{(2;1)}
\deg\mathbf y_{c_1}^{(2;2)}}
\operatorname{Op}
(\mathbf y_{c_1}^{(2;1)})
\otimes
\operatorname{Op}
(\mathbf y_{c_1}^{(2;2)}).
\end{aligned}
\]
第二个等号使用了分次余交换性。
因此，在式~\eqref{eq:7-4-35}
的第二项中令
$\mathbf y=\operatorname{Op}(\mathbf y')$，
它便成为
\[
\begin{aligned}
&\sum_{c_1\in C_1,\,c_2\in C_2}
(-1)^{\substack{
\deg\mathbf y_{c_1}^{\prime(2;2)}
\deg'\mathbf x_{c_2}^{(3;1)}
+\deg'\mathbf x_{c_2}^{(3;1)}\\
+\deg\mathbf y_{c_1}^{\prime(2;1)}}}\\
&\qquad
\mathfrak q\Bigl(
\mathbf y_{c_1}^{\prime(2;1)}
\otimes\bigl(
\mathbf x_{c_2}^{(3;1)}
\otimes
\mathfrak q(
\mathbf y_{c_1}^{\prime(2;2)}
\otimes\mathbf x_{c_2}^{(3;2)})
\otimes\mathbf x_{c_2}^{(3;3)}
\bigr)\Bigr).
\end{aligned}
\]
故式~\eqref{eq:7-4-35}
蕴涵~\eqref{eq:3-8-33}。
于是得到：

\begin{corollary}{7.4.40}\label{cor:7-4-40}
下列两类对象之间存在双射。
\begin{enumerate}[label=\textup{(\arabic*)},leftmargin=2.8em]
\item 严格且有隙的滤过
$L_\infty$ 同态
\[
\widehat{\mathfrak q}^{\circ}:
C(M;\Lambda_{0,\mathrm{nov}})[1]
\longrightarrow
\operatorname{Der}B
\bigl(C(L;\Lambda_{0,\mathrm{nov}})[1]\bigr).
\]

\item 定理~\ref{thm:3-8-32}
中的一族算子
$\mathfrak q_{\ell,k,\beta}$，
满足~\eqref{eq:3-8-33}
和~\eqref{eq:3-8-34-2}。
\end{enumerate}
\end{corollary}

这里在
$C(M;\Lambda_{0,\mathrm{nov}})[1]$
上赋予平凡的滤过
$L_\infty$ 结构。
